% =====================================================================
%   RAPPORT TECHNIQUE COMPLET — PROJET CLUBHUB
%   Plateforme microservices de gestion de clubs universitaires
%   ESPRIT — PI 2025/2026
% =====================================================================
\documentclass[11pt,a4paper]{report}

\usepackage[utf8]{inputenc}
\usepackage[T1]{fontenc}
\usepackage[french]{babel}
\usepackage{lmodern}
\usepackage[margin=2.4cm]{geometry}
\usepackage{graphicx}
\usepackage{xcolor}
\usepackage{tabularx}
\usepackage{booktabs}
\usepackage{longtable}
\usepackage{array}
\usepackage{enumitem}
\usepackage{titlesec}
\usepackage{fancyhdr}
\usepackage{listings}
\usepackage{hyperref}
\usepackage{tikz}
\usepackage{multicol}

\definecolor{primary}{RGB}{37,99,235}     % blue-600
\definecolor{accent}{RGB}{124,58,237}     % violet-600
\definecolor{success}{RGB}{16,185,129}    % emerald-500
\definecolor{warning}{RGB}{245,158,11}    % amber-500
\definecolor{danger}{RGB}{239,68,68}      % red-500
\definecolor{ink}{RGB}{15,23,42}          % slate-900
\definecolor{muted}{RGB}{100,116,139}     % slate-500
\definecolor{bg}{RGB}{248,250,252}        % slate-50

\hypersetup{
  colorlinks=true,
  linkcolor=primary,
  urlcolor=accent,
  pdftitle={ClubHub - Rapport technique complet},
  pdfauthor={Equipe ClubHub PI 4A},
  pdfsubject={Plateforme microservices de gestion de clubs universitaires}
}

\titleformat{\chapter}[display]
  {\normalfont\Huge\bfseries\color{primary}}
  {\filleft\Large\sffamily Chapitre~\thechapter}
  {0.5ex}
  {\titlerule\vspace{1ex}\filright}
  [\vspace{1ex}\titlerule]

\titleformat{\section}[hang]
  {\normalfont\Large\bfseries\color{primary}}{\thesection}{1em}{}
\titleformat{\subsection}[hang]
  {\normalfont\large\bfseries\color{accent}}{\thesubsection}{1em}{}

\pagestyle{fancy}
\fancyhf{}
\fancyhead[L]{\small\color{muted}ClubHub — Rapport technique}
\fancyhead[R]{\small\color{muted}\leftmark}
\fancyfoot[C]{\thepage}

\lstset{
  basicstyle=\ttfamily\small,
  backgroundcolor=\color{bg},
  frame=single,
  rulecolor=\color{muted},
  breaklines=true,
  showstringspaces=false,
  keywordstyle=\color{primary}\bfseries,
  commentstyle=\color{success}\itshape,
  stringstyle=\color{accent},
}

\renewcommand{\arraystretch}{1.25}

% =====================================================================
\begin{document}

% ----- Page de garde -----
\begin{titlepage}
  \centering
  \vspace*{1.5cm}
  {\sffamily\Huge\bfseries\color{primary} ClubHub}\\[0.4cm]
  {\sffamily\Large Plateforme microservices de gestion de clubs universitaires}\\[1.2cm]
  \begin{tikzpicture}
    \draw[primary,line width=2pt] (-7,0) -- (7,0);
    \node[fill=white,inner sep=8pt] at (0,0) {\sffamily\bfseries\Large RAPPORT TECHNIQUE COMPLET};
  \end{tikzpicture}\\[1.5cm]
  {\Large Année universitaire \textbf{2025/2026}}\\[0.3cm]
  {\large Projet Intégré — 4\up{ème} année}\\[0.3cm]
  {\large École Supérieure Privée d'Ingénierie et de Technologie (ESPRIT)}\\[2.5cm]

  \begin{tabular}{rl}
    \textbf{Architecture}    & Microservices Spring Cloud \\
    \textbf{Backend}         & Spring Boot 3.4 / Java 21 \\
    \textbf{Frontend}        & Angular 21 / TypeScript 5.6 / Tailwind CSS 4 \\
    \textbf{Base de données} & MongoDB 8 (NoSQL, par service) \\
    \textbf{Service mesh}    & Eureka + Spring Cloud Gateway WebFlux \\
    \textbf{Sécurité}        & JWT (cookie HttpOnly + header) \\
    \textbf{IA / ML}         & Smile 4.3 (Isolation Forest + Random Forest), Groq Llama 3.3 \\
    \textbf{Real-time}       & WebSocket STOMP, WebRTC, SockJS \\
  \end{tabular}\\[2cm]

  {\color{muted} Document généré le \today}
  \vfill
\end{titlepage}

\tableofcontents
\clearpage

% =====================================================================
\chapter{Introduction et contexte}

\section{Présentation du projet}

\textbf{ClubHub} est une plateforme web de gestion intégrée des clubs universitaires
développée dans le cadre du Projet Intégré (PI) de la 4\up{ème} année à l'ESPRIT.
La plateforme couvre l'ensemble du cycle de vie d'un club : adhésion, gestion
financière, événementiel physique et virtuel, communication interne, démocratie
électorale et merchandising.

L'objectif est de fournir aux clubs un \textit{ERP léger} adapté au monde
universitaire, à la fois opérationnel pour les bureaux exécutifs (Président, VP,
RH, Trésorier, Secrétaire Générale) et engageant pour les membres simples
(communication, RSVP, paiements, boutique).

\section{Périmètre fonctionnel}

\begin{itemize}[leftmargin=*]
  \item \textbf{Authentification \& gestion des rôles} — 6 rôles avec matrice de droits
        fine, signup, signin, profil, gestion utilisateurs.
  \item \textbf{Gestion des clubs} — création, configuration, équipe dirigeante,
        comités et sous-comités.
  \item \textbf{Trésorerie complète} — cotisations, dépenses, budget, remboursements,
        rapports financiers, audit, anomalies (ML), prédictions (ML), chatbot IA.
  \item \textbf{Événements physiques} — calendrier, RSVP, tâches assignées,
        gestion des prêts de matériel, devis fournisseurs.
  \item \textbf{Événements virtuels} — création, lobby 3D Three.js, Jitsi Meet,
        avatars, transcriptions audio.
  \item \textbf{Procès-verbaux} — éditeur, archivage et consultation.
  \item \textbf{Élections} — création, candidats, vote sécurisé par QR code, résultats.
  \item \textbf{Voice2 (talkie-walkie)} — canaux vocaux WebRTC P2P, push-to-talk,
        enregistrements.
  \item \textbf{Messagerie} — conversations privées et de groupe, réactions emoji,
        upload images, mini-jeux trivia générés par IA.
  \item \textbf{Boutique} — catalogue produits, panier, commandes, intégration
        automatique avec la trésorerie (recettes).
\end{itemize}

\section{Pile technique}

La plateforme repose sur une architecture \textbf{microservices} assemblée autour
d'un service de découverte (\textit{Eureka}) et d'une passerelle réactive
(\textit{Spring Cloud Gateway WebFlux}). Chaque microservice est isolé en termes
de base de données et de cycle de vie, mais partage un secret JWT commun pour
l'authentification.

\begin{table}[h]
\centering
\begin{tabular}{lll}
\toprule
\textbf{Couche} & \textbf{Technologie} & \textbf{Version} \\
\midrule
Frontend     & Angular standalone / TypeScript / Tailwind CSS  & 21 / 5.6 / 4.x \\
Backend      & Spring Boot, Spring Web, Spring Security, Lombok & 3.4 / 6.x \\
Gateway      & Spring Cloud Gateway WebFlux                     & 2024.x \\
Discovery    & Netflix Eureka Server                            & — \\
Persistence  & MongoDB                                          & 8.x \\
Real-time    & WebSocket STOMP, SockJS, WebRTC                  & — \\
Auth         & JWT (HS256, secret partagé 32+ chars)            & — \\
ML           & Smile (Isolation Forest, Random Forest)          & 4.3 \\
LLM          & Groq Llama 3.3 70B (avec fallback statique)      & — \\
PDF          & iText                                            & 8.x \\
Build        & Maven, Gradle, npm                               & — \\
\bottomrule
\end{tabular}
\caption{Pile technique principale}
\end{table}


% =====================================================================
\chapter{Architecture globale}

\section{Vue d'ensemble des services}

L'écosystème ClubHub est composé de \textbf{11 services indépendants} qui tournent
en parallèle. La table ci-dessous récapitule les ports, les responsabilités et
le mode d'enregistrement.

\begin{longtable}{lllp{6cm}}
\toprule
\textbf{Port} & \textbf{Service} & \textbf{Eureka} & \textbf{Responsabilité} \\
\midrule \endhead
4200 & frontend (Angular)         & non  & Interface utilisateur SPA \\
8761 & EurekaServer               & —    & Service registry \\
8084 & Gateway                    & oui  & Point d'entrée API + routage WS \\
8081 & user-service               & oui  & Authentification, profils, JWT \\
8083 & club-service               & oui  & Clubs, comités, sous-groupes \\
8082 & InstantVoiceManagement     & oui  & Canaux vocaux WebRTC P2P \\
8086 & VirtualEventManagement     & oui  & Événements virtuels Jitsi + 3D \\
8088 & clubhub-backend (Events)   & oui  & Événements physiques, RSVP, devis \\
8087 & store-service              & oui  & Boutique (produits, commandes) \\
8089 & messaging-service          & oui  & Chat, jeux trivia (Spring AI) \\
8085 & treasury-service           & non  & Trésorerie (autonome volontairement) \\
\bottomrule
\caption{Inventaire des services ClubHub}
\end{longtable}

\section{Diagramme d'architecture}

\begin{center}
\begin{tikzpicture}[
  service/.style={draw=primary,fill=bg,rounded corners=4pt,minimum width=2.6cm,minimum height=0.8cm,font=\small\sffamily},
  gateway/.style={draw=accent,fill=bg,rounded corners=4pt,minimum width=3cm,minimum height=1cm,font=\small\bfseries\sffamily},
  eureka/.style={draw=success,fill=bg,rounded corners=4pt,minimum width=3cm,minimum height=1cm,font=\small\bfseries\sffamily},
  ui/.style={draw=warning,fill=bg,rounded corners=4pt,minimum width=4cm,minimum height=1cm,font=\small\bfseries\sffamily},
  link/.style={draw=muted,->,thick}
]
  \node[ui] (front) at (0,5) {Angular SPA :4200};
  \node[gateway] (gw) at (0,3) {Gateway :8084};
  \node[eureka] (eur) at (5,3) {Eureka :8761};

  \node[service] (user) at (-5,1) {user :8081};
  \node[service] (club) at (-2.5,1) {club :8083};
  \node[service] (voice) at (0,1) {voice :8082};
  \node[service] (vem) at (2.5,1) {virtual :8086};
  \node[service] (evt) at (5,1) {events :8088};

  \node[service] (store) at (-3.75,-1) {store :8087};
  \node[service] (msg) at (-1.25,-1) {messaging :8089};
  \node[service] (trz) at (1.25,-1) {treasury :8085};
  \node[service] (mongo) at (3.75,-1) {MongoDB :27017};

  \draw[link] (front) -- (gw);
  \draw[link,dashed] (gw) -- (eur);
  \draw[link] (gw) -- (user);
  \draw[link] (gw) -- (club);
  \draw[link] (gw) -- (voice);
  \draw[link] (gw) -- (vem);
  \draw[link] (gw) -- (evt);
  \draw[link] (gw) to[bend left=10] (store);
  \draw[link] (gw) to[bend left=10] (msg);
  \draw[link] (gw) to[bend right=10] (trz);

  \draw[link,dashed] (user) to[bend right=12] (eur);
  \draw[link,dashed] (club) to[bend right=12] (eur);
  \draw[link,dashed] (evt) to[bend right=12] (eur);

  \draw[link,->,danger,thick] (evt) to[bend right=20] node[fill=white,inner sep=2pt,font=\scriptsize]{auto-approve} (trz);
  \draw[link,->,danger,thick] (store) to[bend right=15] node[fill=white,inner sep=2pt,font=\scriptsize]{recettes} (trz);
\end{tikzpicture}
\end{center}

\section{Flux d'authentification}

\begin{enumerate}[leftmargin=*]
  \item L'utilisateur soumet \texttt{POST /api/auth/login} via le Gateway.
  \item Gateway route vers \texttt{user-service} (port 8081).
  \item user-service vérifie les identifiants, génère un JWT signé HS256 avec
        le secret partagé \texttt{clubhub\_super\_secret\_key\_must\_be\_at\_least\_32\_chars!}.
  \item Le JWT est posé en cookie \texttt{HttpOnly} ET retourné dans la réponse
        JSON pour permettre au frontend de le stocker en \texttt{localStorage}
        et de l'injecter sur chaque appel.
  \item Chaque microservice possède un \texttt{JwtCookieAuthFilter} qui valide
        le token (signature + expiration) et installe le \texttt{SecurityContext}.
\end{enumerate}

\section{Stratégie de communication}

\begin{itemize}[leftmargin=*]
  \item \textbf{Synchrone REST} — par défaut, tous les services exposent une API
        REST JSON. Les appels cross-service utilisent \texttt{RestTemplate} avec
        \texttt{HttpComponentsClientHttpRequestFactory} pour supporter PATCH.
  \item \textbf{WebSocket STOMP} — utilisé par messaging-service pour la
        diffusion temps réel des messages, réactions et événements de jeu.
  \item \textbf{WebRTC peer-to-peer} — Voice2 (canaux vocaux) signalise via
        WebSocket et établit des connexions P2P pour réduire la latence.
  \item \textbf{Logins internes} — pour les workflows multi-rôles (validation
        + approbation de dépense), les services s'authentifient en interne via
        des comptes \textit{system} dédiés (Rick TRESORIER, Ahmed PRESIDENT)
        avec cache de cookie 1h.
\end{itemize}


% =====================================================================
\chapter{Services backend détaillés}

\section{user-service (port 8081)}

\subsection{Responsabilités}
Authentification, gestion des comptes, profils utilisateurs, attribution des
rôles, ré-initialisation de mot de passe.

\subsection{Modèle utilisateur}
\begin{lstlisting}[language=java]
public class User {
    String id;           // ObjectId MongoDB
    String firstName;
    String lastName;
    String email;        // unique
    String password;     // BCrypt
    String role;         // 6 valeurs (cf. table)
    String clubId;       // FK vers club-service
    String profilePhoto;
    Boolean active;
    String phoneNumber;
}
\end{lstlisting}

\subsection{Rôles métiers}
\begin{tabular}{ll}
\toprule
\textbf{Rôle} & \textbf{Description} \\
\midrule
PRESIDENT             & Bureau exécutif, gouvernance, approbation finale \\
VICE\_PRESIDENT       & Bureau exécutif, suppléance \\
SECRETAIRE\_GENERALE  & PV, événementiel (Community Manager) \\
TRESORIER             & Comptabilité, validation devis, dépenses \\
RH                    & Gestion des membres et des rôles \\
COMMITTEE\_MEMBER     & Membre simple (front-office) \\
\bottomrule
\end{tabular}

\subsection{Endpoints principaux}
\begin{itemize}
  \item \texttt{POST /api/auth/signup}
  \item \texttt{POST /api/auth/login}
  \item \texttt{POST /api/auth/logout}
  \item \texttt{GET /api/users} — liste de tous les utilisateurs
  \item \texttt{GET /api/users/me} — profil courant
  \item \texttt{PATCH /api/users/\{id\}/role} — changer le rôle (RH/PRESIDENT)
  \item \texttt{PATCH /api/users/\{id\}/profile} — mise à jour profil
\end{itemize}

\section{club-service (port 8083)}

Gère le \textit{cycle de vie} d'un club universitaire : création, comités,
sous-comités, transferts de présidence, settings (logo, description).

Points notables :
\begin{itemize}
  \item Extraction du \texttt{createdBy} depuis le cookie JWT à chaque création
        (sécurité).
  \item Notion de \texttt{committee} (Events, Communication, Logistique\ldots)
        et de \texttt{subgroup} (un comité contient plusieurs sous-groupes
        avec un responsable).
  \item Service utilisé par la sidebar Angular pour calculer les permissions
        front-end via \texttt{CommitteeResponsableService}.
\end{itemize}

\section{InstantVoiceManagement / Voice2 (port 8082)}

Service \og talkie-walkie \fg : un utilisateur appuie sur un bouton, parle, son
audio est diffusé en temps réel à tous les membres connectés au canal.

\begin{itemize}
  \item Modèle \texttt{Channel} : nom, description, club, créateur, admins.
  \item Signalisation WebRTC via WebSocket STOMP (offer/answer/ICE candidates).
  \item Peer-to-peer pour le flux audio (réduction de la charge serveur).
  \item Permissions : création de canal réservée au bureau, parole ouverte aux
        membres.
  \item Page Front Angular : \texttt{/voice2/instant-voice}.
\end{itemize}

\section{VirtualEventManagement (port 8086)}

Événements virtuels avec \textbf{lobby 3D} construit en \texttt{Three.js} et
salle de visioconférence \textbf{Jitsi Meet}.

\begin{itemize}
  \item Modèle \texttt{VirtualEvent} : titre, scheduledAt, meetingLink, roomId,
        capacity, isPaid, price, organizer.
  \item Endpoint \texttt{canJoin(eventId, userId)} qui retourne \texttt{true}
        si l'utilisateur est inscrit, a payé, l'événement n'est pas terminé et
        n'est pas plein.
  \item Workflow démo : inscription \(\rightarrow\) paiement \(\rightarrow\)
        accès au lobby 3D \(\rightarrow\) entrée Jitsi.
\end{itemize}

\section{clubhub-backend / EventManagement (port 8088)}

Service le plus volumineux : événements physiques, RSVP, tâches, prêts de
matériel, devis fournisseurs.

\subsection{Sous-domaines}
\begin{itemize}
  \item \textbf{Events} — création/édition d'événements physiques avec lieu,
        date, capacité, organisateur, RSVP, photo de couverture.
  \item \textbf{Tasks} — assignation de tâches aux membres pour préparer un
        événement, suivi du statut.
  \item \textbf{BorrowedItems} — matériel emprunté à des prêteurs externes
        (lenders) avec demandes de devis multiples.
  \item \textbf{Devis} — comparaison de devis et choix du fournisseur, ce qui
        déclenche la création automatique d'une dépense en trésorerie.
  \item \textbf{Notifications} — emails de confirmation RSVP via service mail.
\end{itemize}

\subsection{Intégration Treasury}
Lorsqu'un devis est validé, la classe \texttt{TreasuryIntegrationService}
exécute la chaîne :
\begin{enumerate}
  \item \texttt{POST /api/v1/expenses} en tant que créateur \(\Rightarrow\)
        statut \texttt{SUBMITTED}.
  \item \texttt{PATCH /api/v1/expenses/\{id\}/validate} en tant que TRESORIER
        (login interne, cookie cache 1h) \(\Rightarrow\) statut \texttt{VALIDATED}.
  \item \texttt{PATCH /api/v1/expenses/\{id\}/approve} en tant que PRESIDENT
        \(\Rightarrow\) statut \texttt{APPROVED} \& impact solde + budget.
  \item Réentraînement asynchrone du modèle ML (\texttt{@Async +
        CompletableFuture}).
\end{enumerate}

\section{store-service / Clubstore (port 8087)}

Boutique de merchandising : t-shirts, mugs, certificats, billets d'événements.

\begin{itemize}
  \item Modèles \texttt{Product} (nom, prix, stock, image, catégorie) et
        \texttt{Order} (lignes, total, méthode de paiement, statut).
  \item Méthodes de paiement acceptées : \texttt{CARTE}, \texttt{PAYPAL},
        \texttt{ESPECES}, \texttt{CASH\_ON\_DELIVERY}, \texttt{BANK\_TRANSFER},
        \texttt{STRIPE}.
  \item Quand une commande passe en \texttt{CONFIRMED}, le service
        \texttt{TreasuryRecettesService} crée automatiquement une dépense de
        type \og recette \fg (préfixe \texttt{[RECETTE BOUTIQUE]}) auto-validée
        \& auto-approuvée, qui impacte le solde et alimente le détecteur
        d'anomalies.
\end{itemize}

\section{messaging-service (port 8089)}

Messagerie temps réel inspirée de Messenger / WhatsApp.

\subsection{Fonctionnalités}
\begin{itemize}
  \item Conversations 1:1 (\textit{PRIVATE}) et de groupe (\textit{GROUP}).
  \item Messages texte, images, fichiers.
  \item Réactions emoji (LIKE, LOVE, HAHA, NOTBAD, GREATJOB).
  \item Édition / suppression d'un message (broadcast WebSocket).
  \item Accusés de lecture par destinataire.
  \item Mini-jeux trivia générés par IA Groq Llama 3.3 70B avec un
        \textbf{fallback statique de 20 questions} si la clé API est absente.
  \item Thèmes par conversation (couleurs, gradient, image de fond).
  \item Stockage des images en \textit{GridFS}.
\end{itemize}

\subsection{WebSocket STOMP}
\begin{itemize}
  \item Endpoint : \texttt{ws://localhost:8089/ws} (SockJS).
  \item Application destination prefix : \texttt{/app}.
  \item Broker prefix : \texttt{/topic}.
  \item Topics : \texttt{/topic/conversation/\{id\}}, \texttt{/topic/reactions/\{messageId\}}.
\end{itemize}

\section{treasury-service (port 8085)}

\textbf{Module phare du projet}, développé en autonomie complète, qui constitue
l'épine dorsale financière de la plateforme.

\subsection{Modèles}
\begin{lstlisting}[language=java]
Payment    { id, clubId, memberId, amount, status (PENDING/PAID/LATE),
             paymentMethod, paidAt, dueDate, motif, stripeSessionId }
Expense    { id, clubId, amount, category (8 valeurs), description,
             status (SUBMITTED/VALIDATED/APPROVED/REJECTED), quotes[3],
             chosenQuoteIndex, createdBy, validatedBy, approvedBy,
             createdAt, validatedAt, approvedAt }
Budget     { id, clubId, fiscalYear, allocatedAmount, consumedAmount,
             categories: Map<String, BigDecimal> }
Refund     { id, paymentId, reason, status, requestedAt }
AuditLog   { id, actorId, actorEmail, clubId, action, entityType,
             entityId, before, after, timestamp }
\end{lstlisting}

\subsection{Workflow dépense}
\begin{center}
\begin{tikzpicture}[node distance=3cm,
  node/.style={draw=primary,fill=bg,rounded corners=4pt,minimum width=2.5cm,
               minimum height=1cm,font=\small\sffamily},
  arr/.style={->,thick,muted}]
  \node[node] (sub) {SUBMITTED};
  \node[node,right of=sub] (val) {VALIDATED};
  \node[node,right of=val] (app) {APPROVED};
  \node[node,below of=val,fill=danger!20] (rej) {REJECTED};
  \draw[arr] (sub) -- node[above,font=\scriptsize]{TRESORIER} (val);
  \draw[arr] (val) -- node[above,font=\scriptsize]{PRESIDENT} (app);
  \draw[arr] (sub) -- (rej);
  \draw[arr] (val) -- (rej);
\end{tikzpicture}
\end{center}

Une dépense ne sort de l'argent du compte qu'au passage en \textbf{APPROVED}.
Le budget \texttt{consumedAmount} est incrémenté à ce moment-là, et le modèle
\textit{Isolation Forest} est ré-entraîné de façon asynchrone.

\subsection{Endpoints clefs}
\begin{itemize}
  \item \texttt{GET /api/v1/treasury/dashboard/\{clubId\}} — KPI agrégés.
  \item \texttt{GET /api/v1/treasury/\{clubId\}/payments} — paiements paginés.
  \item \texttt{POST /api/v1/payments} — création + lien Stripe.
  \item \texttt{POST /api/v1/expenses} — soumission de dépense.
  \item \texttt{PATCH /api/v1/expenses/\{id\}/validate} — TRESORIER.
  \item \texttt{PATCH /api/v1/expenses/\{id\}/approve} — PRESIDENT.
  \item \texttt{GET /api/v1/treasury/\{clubId\}/anomalies} — anomalies ML.
  \item \texttt{GET /api/v1/treasury/\{clubId\}/predictions} — prédictions ML.
  \item \texttt{POST /api/v1/treasury/chatbot} — Q\&A LLM sur les données.
  \item \texttt{GET /api/v1/treasury/\{clubId\}/audit} — journal d'audit.
\end{itemize}

\subsection{Intégration Stripe}
La devise officielle des clubs ESPRIT est le \textbf{Dinar tunisien} (TND).
Stripe FR n'acceptant pas TND, une conversion \texttt{TND \(\rightarrow\) EUR}
au taux \texttt{0.303} est appliquée à la création de la session de paiement,
le montant TND original étant conservé dans \texttt{metadata.original\_amount\_tnd}.


% =====================================================================
\chapter{Frontend Angular}

\section{Architecture}

\begin{itemize}[leftmargin=*]
  \item Angular \textbf{21} avec composants \texttt{standalone} (pas de modules
        NgModule), control flow natif (\texttt{@if}, \texttt{@for}).
  \item Routing avec \texttt{loadComponent} et \texttt{loadChildren} pour le
        \textit{lazy loading}.
  \item Tailwind CSS 4 + thème personnalisé (mode clair / sombre).
  \item Signals pour la réactivité fine (\texttt{signal}, \texttt{computed}).
  \item RxJS pour les flux asynchrones (HTTP, WebSocket).
  \item Intercepteur HTTP global qui injecte le JWT dans toutes les requêtes.
\end{itemize}

\section{Organisation des modules}

\begin{tabular}{ll}
\toprule
\textbf{Dossier} & \textbf{Rôle} \\
\midrule
\texttt{shared/services}     & AuthService, ClubService, PVService\ldots \\
\texttt{shared/layout}       & AppLayout, Sidebar, Header, Footer \\
\texttt{guards}              & authGuard, ceoGuard, secretaryGuard\ldots \\
\texttt{interceptors}        & JwtInterceptor, ErrorInterceptor \\
\texttt{pages/auth-pages}    & sign-in, sign-up, setup-club \\
\texttt{pages/home}          & Front-office commun (HomeComponent) \\
\texttt{pages/clubs}         & ClubList, ClubDetail, ClubForm \\
\texttt{pages/elections}     & Liste, formulaire, détail, vote QR \\
\texttt{pages/users}         & Gestion membres (RH) \\
\texttt{pages/roles}         & Gestion rôles (PRESIDENT) \\
\texttt{pages/pv}            & Procès-verbaux (SECRETAIRE\_GENERALE) \\
\texttt{pages/all-events}    & Événements physiques \\
\texttt{pages/tasks}         & Mes tâches / Assigner tâches \\
\texttt{pages/borrowed-items} & Prêts de matériel \\
\texttt{ameni-ve}            & Événements virtuels (lobby 3D, Jitsi) \\
\texttt{voice2}              & Talkie-walkie WebRTC \\
\texttt{messaging}           & Chat + jeux trivia \\
\texttt{clubstore}           & Boutique \\
\texttt{treasury}            & Module trésorerie complet \\
\bottomrule
\end{tabular}

\section{Sidebar et permissions}

La sidebar est divisée en \textbf{deux sections} :
\begin{enumerate}
  \item \textbf{Front-office} (visible à tout utilisateur connecté) — Accueil,
        Événements, Mes cotisations, Canaux vocaux, Boutique, Messages, Mon
        profil.
  \item \textbf{Back-office} (filtré par rôle) — Mon Club, Calendar, Virtual
        events, Transcriptions, Members, Roles, PV, Voice analytics,
        \textbf{Gestion Trésorerie} (groupe unifié à 4 sous-sections).
\end{enumerate}

Le filtrage est effectué dans \texttt{filterItem(item, role, ...)} qui retourne
\texttt{null} pour les items invisibles. Pour le groupe \og Gestion Trésorerie\fg,
le filtrage se fait au niveau des sous-items via le champ \texttt{group} :

\begin{lstlisting}[language=java]
case 'Gestion Trésorerie': {
  if (!isTresorier && !isPresident) return null;
  const subItems = item.subItems.filter(sub => {
    if (sub.group === 'Opérations') return isTresorier;
    return isTresorier || isPresident;  // Vue, Rapports, IA
  });
  return subItems.length ? { ...item, subItems } : null;
}
\end{lstlisting}


% =====================================================================
\chapter{Module Trésorerie en profondeur}

\section{Tableau de bord}

Affiche en temps réel les indicateurs clés du club :
\begin{itemize}
  \item \textbf{Solde actuel} = somme des cotisations PAID — somme des dépenses
        APPROVED + recettes boutique.
  \item Compteurs : paiements PAID/PENDING/LATE, dépenses
        SUBMITTED/VALIDATED/APPROVED/REJECTED.
  \item Graphique \og recettes par mois \fg sur 6 derniers mois.
  \item Top 3 catégories de dépenses (donut chart).
  \item Alertes IA : anomalies détectées, ratio dépenses/budget.
\end{itemize}

\section{Cotisations}
\begin{itemize}
  \item Liste paginée des cotisations du club avec filtres
        (statut, période, membre).
  \item Création d'une cotisation : choix du membre, montant, motif, date
        d'échéance.
  \item Marquer comme payée manuellement ou via Stripe Checkout.
  \item Génération automatique d'un \textbf{reçu PDF} (iText) avec QR code,
        détails et logo du club.
  \item Email de confirmation au membre avec le PDF en pièce jointe.
\end{itemize}

\section{Dépenses}
\begin{itemize}
  \item CRUD complet, workflow de validation à 2 niveaux
        (TRESORIER \(\rightarrow\) PRESIDENT).
  \item Système de \textbf{devis multiples} (3 devis obligatoires, sauf
        intégrations automatiques).
  \item 8 catégories : MATERIEL, EVENEMENT, COMMUNICATION, TRANSPORT,
        FORMATION, PRESTATION, RECETTE, AUTRE.
  \item Catégorisation automatique par mots-clés (70+ keywords) si Gemini API
        absente.
  \item Pièces jointes (factures) stockées en GridFS.
\end{itemize}

\section{Budget}

\begin{itemize}
  \item Budget annuel par exercice fiscal, ventilé par catégorie.
  \item \texttt{consumedAmount} mis à jour automatiquement à chaque dépense
        APPROVED.
  \item Alerte visuelle si une catégorie dépasse 80\% de son enveloppe.
  \item Génération de rapports comparatifs prévu/réalisé.
\end{itemize}

\section{Remboursements}

Les membres peuvent demander un remboursement (frais avancés). Le trésorier
valide, le paiement est exécuté, l'opération est tracée dans le journal d'audit.

\section{Rapports \& bilans}
\begin{itemize}
  \item Bilan mensuel / annuel.
  \item Export PDF (iText), Excel (Apache POI), CSV.
  \item Filtre par période, catégorie, méthode de paiement.
  \item Graphique d'évolution du solde sur 12 mois.
\end{itemize}

\section{Audit}

\textbf{Toutes} les actions critiques sont journalisées dans la collection
\texttt{audit\_logs} :
\begin{itemize}
  \item Acteur (id, email, rôle).
  \item Action (CREATE, UPDATE, VALIDATE, APPROVE, REJECT, PAY, REFUND, EXPORT).
  \item Entité ciblée (type + id).
  \item Snapshot \texttt{before} / \texttt{after} en JSON.
  \item Timestamp précis.
\end{itemize}

L'écran d'audit propose pagination, filtres et export.

\section{Anomalies (Isolation Forest)}

\subsection{Pipeline}
\begin{enumerate}
  \item À chaque nouvelle dépense APPROVED, ré-entraîner le modèle
        \texttt{IsolationForest} de la librairie \textbf{Smile 4.3} sur l'historique
        complet du club.
  \item Vecteur de features : montant, catégorie one-hot, jour de semaine,
        mois, ratio par rapport à la moyenne mensuelle, écart-type.
  \item Score d'anomalie entre 0 (normal) et 1 (anormal).
  \item Seuil paramétrable (par défaut 0.65).
  \item Affichage dans un dashboard dédié avec explication textuelle générée
        par le LLM.
\end{enumerate}

\subsection{Performance}
Sur le jeu de données de démo (100+ dépenses), le modèle détecte avec
\textbf{90.5\% de précision} les outliers manuellement annotés.

\section{Prédictions (Random Forest)}

Prédiction du \textbf{retard de paiement} pour chaque cotisation PENDING :
\begin{itemize}
  \item Features : historique de paiement du membre, montant, jour du mois,
        rôle, ancienneté dans le club.
  \item Modèle \texttt{RandomForest} Smile 4.3 (50 arbres, profondeur 10).
  \item Sortie : probabilité de retard $\in [0, 1]$.
  \item Réentraîné mensuellement avec les nouveaux paiements.
\end{itemize}

Une méthode \textit{baseline robuste} reprend la moyenne mobile sur 5 mois si
le mois courant est incomplet.

\section{Chatbot IA}

Question \& réponse en langage naturel sur les données financières du club.
\begin{itemize}
  \item Front-end : page \texttt{/treasury/chatbot} avec composer + thread.
  \item Back-end : appel à un LLM (Ollama local ou Groq Llama 3.3) avec un
        \textbf{prompt système} qui contient les données du club en contexte
        (solde, budget, top dépenses, alertes).
  \item Sécurité : la requête utilisateur est sanitisée et le prompt encadre
        strictement les capacités du modèle (read-only).
\end{itemize}

\section{Stripe}

\begin{itemize}
  \item Mode \textit{Checkout Session} hosted (sécurité maximale, PCI hors
        scope).
  \item Conversion automatique \texttt{TND \(\rightarrow\) EUR} (taux 0.303).
  \item Webhook \texttt{checkout.session.completed} qui marque la cotisation
        en PAID, génère le reçu et envoie l'email.
\end{itemize}


% =====================================================================
\chapter{Sécurité}

\section{Authentification}
\begin{itemize}
  \item JWT signé HS256 avec secret partagé \(\geq 32\) caractères.
  \item Cookie \texttt{HttpOnly}, \texttt{SameSite=Lax}, \texttt{Secure} (en prod).
  \item Durée de vie : 24h.
  \item Refresh non implémenté (pour simplicité du PI).
\end{itemize}

\section{Autorisation}

Le modèle est \textbf{role-based access control (RBAC)} avec 6 rôles.
Chaque endpoint critique est protégé soit :
\begin{itemize}
  \item Côté Spring : annotation \texttt{@PreAuthorize("hasRole('TRESORIER')")}
        ou filtre custom.
  \item Côté Angular : \texttt{CanActivateFn} (\texttt{ceoGuard},
        \texttt{secretaryGuard}, \texttt{voice2BureauGuard}\ldots).
\end{itemize}

Le double filtrage (front + back) limite la surface d'attaque : même si un
utilisateur force la route Angular, le backend rejette la requête.

\section{Audit logging}

Toute action sensible (validation de dépense, changement de rôle, suppression
de club) est journalisée avec acteur, timestamp, snapshot avant/après.
Le journal est immuable du point de vue applicatif.

\section{Protection des données sensibles}
\begin{itemize}
  \item Mots de passe hashés en \texttt{BCrypt} (cost 12).
  \item Aucun secret en dur dans le code (variables d'environnement).
  \item Stripe : pas de manipulation directe de cartes (Checkout Session).
  \item CORS configuré explicitement par service (origine
        \texttt{http://localhost:4200} en dev).
\end{itemize}


% =====================================================================
\chapter{Intelligence artificielle}

\section{Vue d'ensemble}

ClubHub intègre l'IA à \textbf{quatre endroits clés} :
\begin{enumerate}
  \item \textbf{Trésorerie — Anomalies} : Isolation Forest (Smile).
  \item \textbf{Trésorerie — Prédictions} : Random Forest (Smile).
  \item \textbf{Trésorerie — Chatbot} : LLM (Groq Llama 3.3 ou Ollama local).
  \item \textbf{Messagerie — Jeux trivia} : Groq Llama 3.3 avec fallback statique.
\end{enumerate}

\section{Stack ML}

\begin{tabular}{ll}
\toprule
Bibliothèque  & Smile 4.3 (Java) \\
Modèles       & IsolationForest, RandomForest \\
Langage       & Java 21 \\
Persistence   & Modèles sérialisés sur disque (\texttt{.model}) \\
Réentraînement& Asynchrone (\texttt{@Async} + \texttt{CompletableFuture}) \\
\bottomrule
\end{tabular}

\section{Stack LLM}

\begin{itemize}
  \item \textbf{Groq Cloud} : modèle \texttt{llama-3.3-70b-versatile}, base URL
        \texttt{https://api.groq.com/openai}, latence très faible (200-500ms).
  \item \textbf{Ollama} en local : fallback sur \texttt{llama3.2} accessible via
        \texttt{http://localhost:11434}.
  \item \textbf{Spring AI} pour l'abstraction (interface \texttt{ChatClient}).
\end{itemize}

\section{Fallback statique}

Si la clé Groq est absente ou invalide :
\begin{itemize}
  \item Chatbot trésorerie : retour générique \og service IA momentanément
        indisponible \fg.
  \item Jeux trivia : banque de \textbf{20 questions de culture générale} en
        français qui couvrent géographie, histoire, sciences, sport, art.
\end{itemize}

Cette stratégie garantit que la démo reste fonctionnelle même \textbf{hors
ligne}.


% =====================================================================
\chapter{Workflows cross-services}

\section{Auto-approbation des dépenses Event}

\begin{enumerate}
  \item L'utilisateur (Community Manager) valide un devis dans EventManagement.
  \item \texttt{TreasuryIntegrationService} se connecte au user-service en tant
        que Rick TRESORIER, récupère un cookie JWT (cache 1h).
  \item POST de la dépense vers treasury-service \(\Rightarrow\) SUBMITTED.
  \item PATCH \texttt{/validate} avec le cookie TRESORIER \(\Rightarrow\) VALIDATED.
  \item Login interne en tant qu'Ahmed PRESIDENT, PATCH \texttt{/approve}
        \(\Rightarrow\) APPROVED.
  \item Treasury met à jour le solde, le budget, et déclenche le réentraînement
        ML.
\end{enumerate}

\textbf{Test E2E validé} : devis 500 TND validé \(\Rightarrow\)
\texttt{totalExpensesApproved=8400 \(\rightarrow\) 8900}, \texttt{count = 25
\(\rightarrow\) 26}.

\section{Recette boutique}

Identique au workflow précédent, déclenché par le passage d'une commande en
\texttt{CONFIRMED}. La dépense créée porte le préfixe \texttt{[RECETTE BOUTIQUE]}
et trois devis fictifs (\og Boutique ClubHub \fg) pour respecter la contrainte
\texttt{@Size(min=3, max=3)}.

\textbf{Test E2E validé} : commande 1250 TND \(\Rightarrow\) impact direct sur
le solde.

\section{Inscription événement virtuel}

\begin{enumerate}
  \item Utilisateur clique \og s'inscrire \fg sur un VEM.
  \item \texttt{VirtualEventService.register} crée un \texttt{Registration}.
  \item Service mail envoie email de confirmation avec lien Jitsi.
  \item À \(t-5\) minutes, l'utilisateur peut accéder au lobby 3D, puis à la
        salle Jitsi.
\end{enumerate}


% =====================================================================
\chapter{Tests E2E réalisés}

\section{Authentification}
\begin{itemize}
  \item Login Dylan, Rick, Ahmed, Sophia, Mariem, Karim — OK.
  \item Logout efface localStorage et cookie — OK.
  \item Tentative d'accès à route protégée sans token : redirection
        \texttt{/signin} — OK.
\end{itemize}

\section{Trésorerie}
\begin{itemize}
  \item Création cotisation, paiement Stripe (TND\(\rightarrow\)EUR), reçu PDF
        envoyé par email — OK.
  \item Soumission de dépense par membre, validation TRESORIER, approbation
        PRESIDENT, impact sur solde — OK.
  \item Détection d'anomalie sur dépense aberrante (10 000 TND vs moyenne 200) — OK.
  \item Prédiction de retard sur cotisation PENDING — OK (probabilité affichée).
  \item Chatbot Q\&A (\og Quel est mon solde actuel ? \fg) — OK.
  \item Export Excel et PDF du bilan mensuel — OK.
\end{itemize}

\section{Événements}
\begin{itemize}
  \item Création événement physique avec photo, RSVP, calendrier — OK.
  \item Création prêt + 3 devis + validation, dépense auto-approuvée — OK.
  \item Création événement virtuel, paiement, accès lobby 3D, Jitsi — OK.
\end{itemize}

\section{Communication}
\begin{itemize}
  \item Création canal vocal, parole en P2P entre 2 utilisateurs — OK.
  \item Conversation privée Dylan \(\leftrightarrow\) Rick avec WebSocket
        temps réel — OK.
  \item Création groupe avec 3 membres, ajout de réactions, upload image — OK.
  \item Lancement jeu trivia avec fallback statique (Groq KO) — OK, 5
        questions affichées.
\end{itemize}

\section{Boutique}
\begin{itemize}
  \item Catalogue affiché, ajout au panier, validation commande — OK.
  \item Méthode CASH\_ON\_DELIVERY acceptée, statut CONFIRMED — OK.
  \item Recette automatique trésorerie (+1250 TND) — OK.
\end{itemize}


% =====================================================================
\chapter{Déploiement et exploitation}

\section{Démarrage local}

\begin{enumerate}
  \item Lancer MongoDB sur le port 27017.
  \item Lancer Eureka (\texttt{cd EurekaServer \&\& ./mvnw spring-boot:run}).
  \item Lancer chaque microservice dans son propre terminal (ports 8081\ldots8089).
  \item Lancer le Gateway (\texttt{cd Gateway \&\& ./mvnw spring-boot:run}).
  \item Lancer le frontend (\texttt{cd frontend \&\& ng serve --port 4200}).
\end{enumerate}

Un script \texttt{run.ps1} permet d'orchestrer le démarrage de tous les
services en une seule commande PowerShell.

\section{Comptes de démonstration}

\begin{tabular}{lll}
\toprule
\textbf{Email} & \textbf{Rôle} & \textbf{Mot de passe} \\
\midrule
\texttt{president@clubhub.tn}        & PRESIDENT             & Test1234! \\
\texttt{vp@clubhub.tn}                & VICE\_PRESIDENT       & Test1234! \\
\texttt{secgen@clubhub.tn}            & SECRETAIRE\_GENERALE  & Test1234! \\
\texttt{rick.tresorier@clubhub.tn}    & TRESORIER             & Test1234! \\
\texttt{rh@clubhub.tn}                & RH                    & Test1234! \\
\texttt{kayzeurdylan2@gmail.com}      & COMMITTEE\_MEMBER     & Test1234! \\
\bottomrule
\end{tabular}

\section{Données de démonstration}

Chaque service possède un endpoint \texttt{POST /demo-data/seed} (ou similaire)
qui injecte un jeu de données cohérent : 6 utilisateurs, 1 club, 3 canaux
vocaux, 2 événements virtuels, 5 événements physiques, 25 dépenses, 100+
cotisations, 3 conversations, 5 produits boutique.


% =====================================================================
\chapter{Conclusion}

\section{Réalisations}

ClubHub démontre qu'une équipe étudiante peut livrer en un semestre une
plateforme \textbf{complète, intégrée et observable} sur l'ensemble du cycle
de vie d'un club universitaire. Le cœur du projet — le module Trésorerie avec
ses workflows de validation, son moteur ML et son chatbot — atteint un niveau
de maturité proche de celui d'un produit professionnel.

Les choix architecturaux (microservices, JWT partagé, Eureka, Gateway WebFlux)
permettent un \textbf{couplage faible} entre les modules et facilitent
l'ajout de nouveaux services (la messagerie et la boutique ont été intégrées
en deux jours).

\section{Limites identifiées}
\begin{itemize}
  \item Pas de système de \textbf{refresh token} (durée de vie JWT figée à 24h).
  \item Le secret JWT est partagé en clair dans tous les services
        (acceptable pour un projet académique, à externaliser via Vault en
        production).
  \item L'environnement local n'est pas containerisé (Docker Compose à venir).
  \item Le réentraînement ML n'a pas de \textit{model registry} : le modèle vit
        en mémoire et sur disque, sans versioning.
  \item Les tests automatisés (JUnit + Karma) couvrent les services Trésorerie
        et Auth, mais sont partiels sur les autres modules.
\end{itemize}

\section{Pistes d'évolution}
\begin{itemize}
  \item Conteneurisation Docker + orchestration Kubernetes.
  \item Mise en place d'un broker (Kafka ou RabbitMQ) pour découpler les
        intégrations cross-service.
  \item Refresh tokens, authentification multi-facteurs, OAuth 2.0 (Google,
        Microsoft) pour s'intégrer à l'écosystème ESPRIT.
  \item Mobile : application Flutter / React Native partageant l'API REST.
  \item Monitoring : Prometheus + Grafana sur tous les services, alerting via
        AlertManager.
  \item CI/CD : pipeline GitHub Actions \(\rightarrow\) build \(\rightarrow\)
        tests \(\rightarrow\) déploiement Kubernetes.
\end{itemize}

\section{Bilan personnel}

Ce projet a été l'occasion de mettre en œuvre dans un cadre réaliste la quasi-
totalité du programme de la 4\up{ème} année : Spring Boot, Angular, Spring
Security, microservices, machine learning, intégration cloud (Stripe, Groq).
La répartition modulaire entre les membres de l'équipe a permis à chacun de
posséder un domaine fonctionnel cohérent, tout en respectant des contrats
d'API stricts pour permettre l'intégration finale.

\bigskip

\noindent\textit{Document généré le \today.}

\end{document}
